%------------------------------------------------------------------------------%
\section{Potências}

%------------------------------------------------------------------------------%
\begin{definicao}{Potência}{def:potencia}
  Seja $a$ um \emph{número real positivo} chamado \emph{base}.
  Definimos para todo $n\in\N$, a \emph{Potência}\index{Potência}
  de $a$ elevado a $n$, como
  \[
    a^n = \underbrace{\,a\,a\,a\,\cdots\, a\,}_{n\text{ vezes}}.
  \]
\end{definicao}
%------------------------------------------------------------------------------%

Em particular,
\[
  a^1 = a.
\]
Segue da Definição~\ref{def:potencia} que para todo
$m, n\in\N$ as seguintes propriedades são válidas
\begin{gather}
  a^m a^n = a^{m+n}                               \label{eq:protencia_1}  \\[2mm]
  \left(a^m\right)^n = a^{mn}                     \label{eq:protencia_2}  \\[2mm]
  \text{Se } b>0, \text{ então } (ab)^m = a^m b^m \label{eq:protencia_3}
\end{gather}

Fixada uma base $a$, queremos estender a Definição~\ref{def:potencia} para o
conjunto dos números reais, ou seja, definir $a^x$ para todo $x\in\R$, de tal
forma que as propriedades acima continuem válidas.

Observe que fixamos uma base $a$, ou seja, $a$ por definição é maior que zero.
Posteriormente falaremos do porquê dessa escolha.

Faremos essa extensão gradualmente.

%------------------------------------------------------------------------------%
\emph{Como definir $a^0$?}

Fazendo $m=0$ e $n=1$ em~\eqref{eq:protencia_1} temos que
\[
  a^0 \, a^1 = a^{0+1} = a^1
\]
Como $a\neq 0$, dividimos ambos os lados da equação por $a=a^1$ e
concluímos que $a^0=1$. Portanto, definimos
\begin{equation}
  \label{eq:potencia_zero}
  a^0 = 1.
\end{equation}
Observe que $0^0$ é uma expressão não definida, ou seja, $0^0$ não é um
número real.

%------------------------------------------------------------------------------%
\emph{Como definir $a^{-m}$ para $m$ inteiro positivo?}

Novamente aplicando a propriedade~\eqref{eq:protencia_1} e utilizando
que $a^0=1$, temos que
\[
  a^{-m}\, a^{m}
  = a^{-m+m}
  =a^0
  =1.
\]
Portanto, definimos
\begin{equation}
  \label{eq:potencia_negativa}
  a^{-m} = \frac{1}{a^m}.
\end{equation}

%------------------------------------------------------------------------------%
\begin{example}
  Avalie a potência \(3^{-2}\).

  \solution
  \[
    3^{-2}
    = \frac{1}{3^2}
    = \frac{1}{3\cdot 3}
    = \frac{1}{9}.
  \]
\end{example}

%------------------------------------------------------------------------------%
\begin{example}
  Avalie a potência \(\left(\frac{2}{3}\right)^{-3}\).

  \solution
  \begin{align*}
    \left(\frac{2}{3}\right)^{-3}
    & = \frac{1}{\left(\dfrac{2}{3}\right)^3}  \\[2mm]
    & = \frac{1}{\dfrac{2^3}{3^3}}             \\[2mm]
    & = \frac{3^3}{2^2}                        \\[2mm]
    & = \left(\frac{3}{2}\right)^{3}           \\[2mm]
    & = \frac{3}{2}\;\frac{3}{2}\;\frac{3}{2}  \\[2mm]
    & = \frac{27}{8}.
  \end{align*}
\end{example}
%------------------------------------------------------------------------------%

Observem que se $a$ e $b$ são números positivos, então
\begin{equation}
  \left(\frac{a}{b}\right)^{-n}=\left(\frac{b}{a}\right)^{n}.
\end{equation}

%------------------------------------------------------------------------------%
\emph{Como definir $a^x$ para $x\in\Q$?}

Inicialmente faremos um exemplo.

%------------------------------------------------------------------------------%
\begin{example}
  Avalie \(a^{\sfrac{n}{2}}\).

  \solution
  Usando a propriedade~\eqref{eq:protencia_2} temos
  \[
    a^{\sfrac{n}{2}}\, a^{\sfrac{n}{2}}
    = \left(a^{\sfrac{n}{2}}\right)^2
    = a^{2\frac{n}{2}}
    = a^n.
  \]
  Como o produto de \(a^{\sfrac{n}{2}}\) por ele mesmo resultou em $a^n$,
  concluímos que
  \[
    a^{n/2} = \sqrt{a^n}.
  \]
\end{example}
%------------------------------------------------------------------------------%

De forma análoga, a propriedade~\eqref{eq:protencia_2} pode ser utilizada para
definir $a^{\sfrac{n}{m}}$ para todo \(\sfrac{n}{m}\in\Q\),
a partir da função
\[
    f(x) = \sqrt[m]{x}.
\]
Como $a^n$ já está definida para todo $n\in\N$
e ao multiplicarmos $a^{\sfrac{n}{m}}$ por si mesmo $m$ vezes, obtemos
\[
  \underbrace{\,a^{\sfrac{n}{m}}\,\cdots\,a^{\sfrac{n}{m}}\,}_{m\text{ vezes}}
  = \left(a^{\sfrac{n}{m}}\right)^m
  = a^{m \frac{n}{m}}
  = a^n,
\]
temos que
\begin{equation}
  a^{\sfrac{n}{m}} = \sqrt[m]{a^n}.
\end{equation}

%------------------------------------------------------------------------------%
\begin{example}
  Avalie a potência \(3^{\sfrac{5}{2}}\).

  \solution
  \[
    3^{\sfrac{5}{2}}
    = \sqrt{3^5}
    = \sqrt{3^4\cdot 3}
    = 9\sqrt{3}.
  \]
\end{example}
%------------------------------------------------------------------------------%

%------------------------------------------------------------------------------%
\begin{example}
  Avalie a potência \(\left(\frac{2}{3}\right)^{-\sfrac{2}{3}}\).

  \solution
  \[
    \left(\frac{2}{3}\right)^{-\sfrac{2}{3}}
    = \left(\frac{3}{2}\right)^{\sfrac{2}{3}}
    = \sqrt[3]{\,\dfrac{9}{4}\,}.
  \]
\end{example}
%------------------------------------------------------------------------------%

%------------------------------------------------------------------------------%
\emph{Como definir $a^x$ para $x\in\mathbb{I}$?}

Fixado $x$ irracional, existe uma sequência de números racionais $x_n$ tal
que os valores dessa sequência ficam cada vez mais próximos do número
irracional $x$. Por exemplo, se
\[
  x
  = \pi
  =3,141\,592\,653\,589\dots,
\]
podemos considerar a sequência dada por
\begin{align*}
  x_1 & = 3          \\
  x_2 & = 3,1        \\
  x_3 & = 3,14       \\
  x_4 & = 3,141      \\
  x_5 & = 3,1415     \\
  x_6 & = 3,14159    \\
  x_7 & = 3,141592   \\
  \vdots
\end{align*}
Nesse caso, dizemos que o limite quando $n$ tende a infinito da sequência
$x_n$ converge para $x$ (estudaremos com mais detalhe a ideia de limite no
próximo módulo, nesse momento, apenas note que quando $n$ é muito grande o
número $x_n$ fica muito próximo do número irracional $x$) e escrevemos
\[
  \lim_{n\to\infty}x_n = x.
\]
Para cada $x_n$, $a^{x_n}$ já foi definido. A sequência $y_n=a^{x_n}$
converge para um valor (ou seja, à medida que $n$ cresce, os valores
$a^{x_n}$ ficam muito próximos de um número real que chamaremos de $a^x$).
Definimos então
\begin{equation}
  a^{x} = \lim_{n\to\infty}a^{x_n}.
\end{equation}

%------------------------------------------------------------------------------%